\section*{Abstract}\label{abstract}

Service composition for moving IoT devices is fundamentally different
from static environments. When services move---pedestrians in a shopping
center or walking through a campus---the composition that worked seconds
ago may fail moments later. This paper presents an Advantage
Actor-Critic (A2C) approach for composing moving IoT services based on
spatio-temporal validity: a service is valid only when it falls within
the consumer's discovery zone at the current time instant.

We build upon the Signal Transmission Reward (STR) model from prior work [1], which derives service quality from Euclidean distance through exponential signal attenuation and Shannon-Hartley capacity. The A2C
agent learns to select valid services that maximize capacity while
avoiding services outside the communication range.

We evaluate on two real GPS trajectory datasets: the ATC shopping center
dataset (185,554 trajectories, 1.7 billion samples from Osaka, Japan)
and the Illinois daily commute dataset (207 trajectories, 357,706
samples). Our implementation compares shared and separate network
architectures. Shared networks converge faster in low-mobility
scenarios (2-5 km/h pedestrian speed), while separate networks achieve
up to 92.3\% valid action selection on real-world pedestrian data.
The A2C approach reduces re-composition frequency by 44\% through
entropy-regularized policy learning.

\textbf{Keywords}: Moving IoT services, service composition, Advantage
Actor-Critic, spatio-temporal constraints, deep reinforcement learning,
Signal Transmission Reward

\section{Introduction}
